\cotefile{157-2}
\scanpage{163}\mspage{108}
\uline{Théorème-Scholie 2} Soit $\unsure{M|\mathcal M}$ une catégorie, munie de deux ensembles de flèches
\[
W,\mathrm{Fib}\subset\mathrm{Fl}(M).
\]

\deleted{On supposera que}\inserted{On fait les hypothèses suivantes} 1°) $W$ contient les \unsure{Iso|Isos}, \deleted{\unsure{et que}} si $f,g$ sont composables, si deux parmi $f,g,gf$ sont dans $W$, la troisième aussi.

2°) \deleted{\unsure{$\mathrm{Fib}$}} $M$ est stable par $\underset{\leftarrow}{\lim}$ finies, $\mathrm{Fib}$ \deleted{\unsure{est}} contient les \unsure{Isomorphismes|isomorphismes}, est stable par composition et par changement de base.
% slanted
\margin{Ces hypothèses \inserted{de stabilité} sur $\mathrm{Fib}$ ne serviront que dans II, et encore seulement (!) dans le cas A), repris in toto \unsure{sous} la forme A$'$ plus bas (p.~\unsure{112})}

\inserted{Pour $S$ dans $M$, $\boxed{\underline{\mathrm{Fib}}\,S}$ désigne la ss-catégorie pleine de $M_{/S}$ formée des $X\longrightarrow S$ qui sont des $\mathrm{Fib}$.}

On distingue deux \deleted{\unsure{conditions pour}} \inserted{versions différentes} d'un axiome \deleted{\unsure{(A)}} de factorisation d'une $f:X\longrightarrow S$ quelconque dans $M$ en $X\xrightarrow{i}\bar X\xrightarrow{\bar f}S$, avec $i\in W$, $\bar f\in\mathrm{Fib}$.

A) On se donne, pour \deleted{\unsure{tout}} $S$ dans $M$, un foncteur
\[
M_{/S}\longrightarrow M_{/S}\qquad X\longmapsto C(X)\quad(=\bar X),
\]
\deleted{avec $C(X)$ \unsure{et}} et un morphisme de foncteurs en $X$
\[
X\xrightarrow{i_X}C(X),
\]
avec
\[
i_X\in W,\qquad C(X)\in\mathrm{Ob}\,\underline{\mathrm{Fib}}\,S.
\]

B) Désignant par
\[
\mathrm{Cof}_0\subset\mathrm{Fl}\,M
\]
l'ens. des flèches de $M$ qui possèdent la propriété de relèvement à droite p.r. à $\mathrm{Fib}$, on suppose ``à la Quillen'' que toute \deleted{\unsure{application}} flèche $f:X\longrightarrow S$ se factorise en
\[
X\xrightarrow{i_{\deleted{\unsure{X}}}}C(\deleted{\unsure{\bar X}}X)\xrightarrow{\bar f}S,
\]
avec
\[
i\in\mathrm{Cofib}_0\cap W,\quad \bar f\in\mathrm{Fib}.
\]
% archive sheet number: 162
\scanpage{164}\mspage{109}
Donnons nous $S$ dans $\unsure{M|\mathcal M}$, considérons
\[
W_S\subset\mathrm{Fl}(M_{/S})
\]
l'ens de flèches induit par $W$ ($f:$
\diagram{\begin{tikzcd}[ampersand replacement=\&,column sep=small,row sep=tiny]
X \arrow[rr] \arrow[dr] \& \& Y \arrow[dl] \\
\& S \&
\end{tikzcd}}
est dans $W_S$ ssi il l'est en oubliant $S$), et désignons par $W_S$ (par abus de notation) l'ens de flèches induit par $W_S$ sur $\underline{\mathrm{Fib}}\,S$, ou plus généralement sur \deleted{\unsure{toute}} sous-catégorie pleine $\unsure{\mathcal F|F}$ de $M_{/S}$, avec
\[
\underline{\mathrm{Fib}}\,S\subset\mathcal F\subset M_{/S}.
\]

\textcircled{\unsure{I|II}} \deleted{\illegible{1}} Le foncteur
\[
(\underline{\mathrm{Fib}}\,S)W_S^{-1}\xrightarrow{\simeq}(M_{/S})W_S^{-1}\tag*{$(*)$}
\]
\inserted{induit par l'inclusion} est une équivalence de catégories, plus généralement pour \deleted{\unsure{toute}} $\mathcal F$ comme dessus
\[
(\underline{\mathrm{Fib}}\,S)W_S^{-1}\longrightarrow\mathcal F W_S^{-1}\longrightarrow(M_{/S})W_S^{-1}
\]
sont des \unsure{équiv.} de catégories. Un foncteur quasi-inverse est obtenu (dans les deux cas A) et B)) à l'aide du foncteur
\[
\gamma:M_{/S}\longrightarrow(\underline{\mathrm{Fib}}\,S)W_S^{-1}
\]
qui est donné sur les objets par
\[
\gamma(X)=C(X),
\]
et qui sur les flèches se définit de façon évidente dans le cas A (\unsure{fonctorialité|fonctorialité ?} de $C(X)$ en $X$), et dans le cas de B se construit ad-hoc, comme la seule loi fonctorielle pour laquelle les hom
\[
i_X:X\longrightarrow C(X)
\]
% archive sheet number: 163
\scanpage{165}\mspage{110}
(interprétés comme flèches dans $(\unsure{M|\mathcal M}_{/S})W_S^{-1}$) soient eux-mêmes fonctoriels en $X$.

On désigne par
\[
\mathrm{Ho}_W^{\mathrm{lc}}(S)\qquad\text{(lc pour ``loc. constant'').}
\]
la catégorie de fractions $(M_{/S})W_S^{-1}$, équivalente à $(\underline{\mathrm{Fib}}\,S)W_S^{-1}$.

\textcircled{II} Donnons nous un localiseur $W_S^{\unsure{0|o|\circ}}$ de $M_{/S}$
\[
W_S^0\subset W_S\subset\mathrm{Fl}(M_{/S}),
\]
avec \unsure{les hypothèses|la hypothèse} formelles similaires à celles de $W$ (cf 1.). On suppose que
\[
W_S^0\cap\mathrm{Fl}(\underline{\mathrm{Fib}}\,S)=W_S\mid\underline{\mathrm{Fib}}\,S.\tag*{3°)}
\]

(Cf pages 94--96 pour une façon générale d'obtenir de tels localiseurs. Il faut bien voir que pour définir des ``types ``d'homotopie'' relatifs sur $S$'', le localiseur $W_S$ est beaucoup trop grossier, \unsure{manifestement|moralement|évidemment} il \inserted{ne} donne qu'une catégorie des types d'hom. rel. localement constants, \unsure{\})} et c'est justement ça qui est montré par \deleted{\unsure{l'}} l'équivalence de \textcircled{I}, où $\underline{\mathrm{Fib}}\,S$ est visualisé comme formé des ``fibrations loc. triviales'' sur $S$.)
On a alors un deuxième localisé
\[
\mathrm{Ho}_{W_S^0}(S)\overset{\mathrm{def}}{=}(M_{/S})(W_S^0)^{-1},
\]
et deux foncteurs
\diagram{\begin{tikzcd}[ampersand replacement=\&]
{(*)} \& \mathrm{Ho}_{W_S^0}(S) \arrow[r,shift left,"\rho"] \& \mathrm{Ho}_W^{\mathrm{lc}}(S) \arrow[l,shift left,"i"]
\end{tikzcd}}
où $\rho$ est le foncteur de localisation \inserted{(strict)} évident, dû à l'inclusion $W_S^0\subset W_S$, et où $i$ est
% archive sheet number: 164
\scanpage{166}\mspage{111}
% archive sheet number: 165
essentiellement le foncteur déduit de l'inclusion
\[
\left(\underline{\mathrm{Fib}}\,S,\underbrace{W_S\mid\underline{\mathrm{Fib}}\,S}_{\shortparallel\,W_S^{\unsure{0|o|\circ}}\mid\underline{\mathrm{Fib}}\,S}\right)\longrightarrow(\unsure{M|\mathcal M}_{/S},W_S^0)
\]
par localisation au \uline{moyen de} $W_S^0$, donc le \deleted{\illegible{1}} foncteur
\[
(\underline{\mathrm{Fib}}\,S)(W_S^0)^{-1}\longrightarrow(M_{/S})(W_S^0)^{-1},\tag*{$(**)$}
\]
compte tenu de l'équivalence $(*)$ de \textcircled{I} ; i.e. $i$ est le composé du quasi-inverse de $(*)$, avec le foncteur $(**)$ précédent. Explicitement, il est défini \unsure{sur|par}
\[
i(X)=\text{Classe de }C(X)\text{ dans }(M_{/S})(W_S^0)^{-1}\,;
\]
en tenant compte du fait que $X\longmapsto C(X)$ transforme $W_S$ en $W_S^0$, grâce à la condition 3°) sur $W_S^0$.

Ceci posé, \uline{sous l'hyp. B} (``Quillen''), on a : dans le couple de foncteurs $(\rho,i)$ de $(*)$ page 110, \uline{$i$ est adjoint à droite de $\rho$}, i.e. on a pour $\xi\in\mathrm{Ob}\,\mathrm{Ho}_{W_S^0}(S)$, $\eta\in\mathrm{Ho}^{\mathrm{lc}}_W(S)$
\[
\mathrm{Hom}_{W_S^0}(\xi,i(\eta))\simeq\mathrm{Hom}_{W_S}(\rho(\xi),\eta),
\]
ce qui s'explicite, pour $X,Y$ dans $\deleted{\unsure{\in}}M_{/S}$, par
\[
\boxed{\mathrm{Hom}_{W_S^0}(X,C(Y))\simeq\mathrm{Hom}_{W_S}(X,Y)}
\]
% slanted
\margin{\uline{NB} le premier membre est indépendant du choix de $W_S^0$ ! (Pourvu seulement que l'hyp $W_S^0\cap\mathrm{Fl}\,\underline{\mathrm{Fib}}\,S=W_S\cap\mathrm{Fl}\,\underline{\mathrm{Fib}}\,S$, en plus de $W_S^0\subset W_S$, soit satisfaite.)}
\unsure{S'ensuit|Il s'ensuit} que \uline{le foncteur $i$ est pleinement fidèle.}

La même formule d'adjonction, et la
\scanpage{167}\mspage{112}
pleine fidélité de $i$, sont valables dans le cas $\deleted{\unsure{B}}A$, quand on renforce l'hypothèse A (relative à un $S$ particulier) en une hypothèse plus restrictive et plus compréhensive, que voici.

A$'$) Le foncteur $C:\unsure{M|\mathcal M}_{/S}\longrightarrow M_{/S}$ et \unsure{le|les} \deleted{\unsure{immersions}} \unsure{morphisme|morphismes} \inserted{\unsure{fonctoriel|fonctoriels}} $i_X:X\longrightarrow C(X)$ sont construits ``à la \unsure{Cartan-Serre}'' à partir d'un foncteur défini sur $M$ \deleted{tout} entier
\[
M\xrightarrow{\mathrm{Ch}}M,
\]
du type ``chemins'', i.e. muni en plus des données suivantes
\[
X\xrightarrow{\alpha_X}\mathrm{Ch}(X)\xrightarrow{p_X=(\sigma_X,\beta_X)}X\times X\tag*{a)}
\]
(structure de source et but des chemins, et de chemins \deleted{\unsure{identiques}}\inserted{triviaux}), et

b) Une composition des chemins, sans exiger ici l'associativité de la composition, mais seulement que les chemins triviaux soient \unsure{neutres|unités} bilatères pour la composition.

De plus, on suppose que pour tout $X$
\[
\left\{\begin{aligned}
&\alpha_X\in W,\quad p_X\in\mathrm{Fib}\\
&(X\longrightarrow e)\in\mathrm{Fib}.
\end{aligned}\right.\tag*{\deleted{\unsure{(3)}}}
\]
% archive sheet number: 166
\scanpage{168}\mspage{113}
\uline{Corollaire} Soient $X$ dans $\unsure{M|\mathcal M}_{/S}$, $Y$ dans $\underline{\mathrm{Fib}}\,S$, alors
\[
\mathrm{Hom}_{W_S^{\unsure{0|o|\circ}}}(X,Y)\xrightarrow{\sim}\mathrm{Hom}_{W_S}(X,Y).
\]

\uline{Dém} Par la formule d'adjonction (p.~111), le deuxième membre est \deleted{\unsure{$\simeq$}} $\mathrm{Hom}_{W_S^0}(X,CY)$. Mais \deleted{\unsure{comme}} $Y$ \inserted{et $C(Y)$} \deleted{\unsure{est}}sont dans $\underline{\mathrm{Fib}}\,S$, où $W_S$ et $W_S^0$ coïncident, $Y\xrightarrow{i_{\unsure{X|Y}}}C(Y)$ est non seulement dans $W_S$, mais dans $W_S^0$, donc $\mathrm{Hom}_{W_S^0}(X,Y)\xrightarrow{\sim}\mathrm{Hom}_{W_S^0}(X,CY)$, d'où \unsure{le|la} résultat \deleted{\illegible{1}}

\figure{\begin{tikzpicture}\draw (0,0) -- (0.52,-0.58);\end{tikzpicture}}

Revenons à $\mathrm{Cat},\mathrm{Cat}_{/S}$. La digression précédente sur $(M,W,\mathrm{Fib})$ a permis de mieux situer ce qui concerne le foncteur $C$, et les relations entre les catégories localisées de
\[
\begin{aligned}
\mathrm{Cat}_{/S}\quad(\text{pour }W_S,W_S^d,W_S^g)\quad&\text{et }\underline{\mathrm{Parf}}\,S\quad(\text{pour }W_S^{\unsure{f}}={}\\
&W_S=W_S^d=W_S^g)
\end{aligned}
\]
ainsi que pour les catégories intermédiaires
\[
\underline{\mathrm{Parf}}\,S\subset\unsure{\mathcal F|F}\subset\mathrm{Cat}_{/S}
\]
pleines de $\mathrm{Cat}_{/S}$ (sur le diagramme p.~76, il y en a cinq autres, savoir $\underline{\unsure{\mathrm{Fibhot}|\mathrm{Fiblat}}}\,S$, $\underline{\mathrm{Liss}}\,S$, $\underline{\mathrm{Liss}}_0\,S$, $\underline{\mathrm{Prop}}\,S$, $\unsure{\mathrm{Prop}_0}\,S$).

Pour ce qui concerne les relations entre les localisés de ces sept catégories, qui se trouvent au centre du diagramme, les énoncés du th.~scholie 1 (p.~\unsure{70|80}--81) ne donnent rien de plus que \unsure{le contenu|les contenus} du th.~scholie 2 (p.~108--113). \deleted{\unsure{En}} Pour les établir, on a donc le choix entre \inserted{les} deux hypothèses A) ou B) dudit théorème. Soit utiliser un foncteur $C_S$ explicite, à l'aide d'un
% archive sheet number: 167
\scanpage{169}\mspage{114}
foncteur $\inserted{X\longmapsto}\underline{\mathrm{Ch}}(X)$ (``chemins dans $X$'') sur $\mathrm{Cat}$, avec les structures de source, but et composition pour les chemins, et des chemins triviaux (agissant comme unités bilatères pour la composition). La construction d'un tel foncteur $C$ n'a pas l'air trivial et semble ne pas être connue -- j'ai peiné pour y arriver ! (\unsure{Avec|Avoir}, en plus, l'associativité de la composition \ldots) Ou bien utiliser, à la Quillen-\deleted{Thomason}, l'existence d'une classe de morphismes dits Q-fibrants, avec des bonnes propriétés de factorisation.
Ce n'est pas non plus évident \unsure{du|de|dans} $\mathrm{Cat}$, et c'est établi par Thomason, par une voie un peu compliquée. Peut-être serai-je amené à revenir dessus -- mais seulement en cas de besoin !

Par contre les \unsure{assertions} concernant les catégories $\underline{\mathrm{Fib}}\,S$, $\underline{\mathrm{Fibsc}}\,S$, ou $\underline{\mathrm{Cofib}}\,S$, $\underline{\mathrm{Cofibsc}}\,S$, et leurs variantes à indices $0$ (donc les huit catégories-\unsure{ailes|astres} du diagramme), et leurs relations à $\mathrm{Cat}_{/S}$ d'une part, à $\underline{\unsure{\mathrm{Fibhot}|\mathrm{Fiblat}}}\,S$ de l'autre (pour les cat. à indice $0$), ne peuvent être contenues dans un topo général sur des \deleted{\unsure{localiseurs}} cat. de modèles $(\unsure{M|\mathcal M},W)$, \inserted{me semble-t-il,} vu la nature très spéciale des quatre catégories désignées ci-dessus, de leur rôle très particulier dans le \deleted{\unsure{théorème}} formalisme catégorique. Or tout au moins les catégories ``scindées'' $\underline{\mathrm{Fibsc}}\,S$, $\underline{\mathrm{Cofibsc}}\,S$, et par suite aussi leurs variantes à indice $0$, jouent un rôle essentiel pour la théorie de $\mathrm{HOT}(S)$, vu que c'est elles qui
% archive sheet number: 168
\scanpage{170}\mspage{115}
\deleted{\unsure{correspondent à la}}\inserted{interviennent directement} quand on applique la recette générale de construction d'un dérivateur à partir d'une catégorie de modèles $(\unsure{M|\mathcal M},W)$. Donc je ne vois pas de simplification substantielle à apporter au th.-scholie 1. Tout au plus pourrait-on élaguer le diagramme des catégories envisagé, en laissant tomber $\underline{\mathrm{Liss}}\,S$, $\underline{\mathrm{Liss}}_0\,S$, $\underline{\mathrm{Prop}}\,S$, $\underline{\mathrm{Prop}}_0\,S$, qui n'y jouent pas un rôle vraiment utile.
Pour ce qui concerne les résultats énoncés, \unsure{ce|en ce} qui les concerne est couvert par ce qui est dit pour les sous-catégories pleines $\unsure{\mathcal F|F}$ de $\mathrm{Cat}_{/S}$ qui contiennent (disons) $\underline{\mathrm{Parf}}\,S$ (ou au choix, celles qui contiennent $\underline{\mathrm{Fib}}\,S$ ou $\underline{\mathrm{Cofib}}\,S$).

\figure{\begin{tikzpicture}\draw (0,0) -- (0.36,-0.66);\end{tikzpicture}}

La description des catégories du type $\mathrm{HOT}(S)$, pour $S$ fixé, me paraît à présent parfaitement bien comprise. Il est temps de faire varier $S$, et de comprendre la variation contravariante en $S$. On se donne donc
\[
S'\xrightarrow{t}S,
\]
et on se propose de décrire les foncteurs images inverses pour les trois types de catégories $\mathrm{HOT}(S)$. Mais nous allons nous borner à présent aux $\mathrm{HOT}_W(S)$ et $\mathrm{HOT}^{\mathrm{lc}}_W(S)$, puisque tout ce qui concerne $\mathrm{HOT}^{\unsure{o|0|\circ}}_W(S)=\mathrm{HOT}_W(S^o)$ \deleted{\unsure{s'en}} se déduit par dualité $S\longmapsto S^o$. Donc il s'agit de décrire
% archive sheet number: 169
\scanpage{171}\mspage{116}
\[
f^*:\mathrm{HOT}_W(S)\longrightarrow\mathrm{HOT}_W(S')\tag*{$(*)$}
\]
et le foncteur induit
\[
f^*:\mathrm{HOT}^{\mathrm{lc}}_W(S)\unsure{\longrightarrow|\text{---}}\mathrm{HOT}^{\mathrm{lc}}_W(S').
\]

Quand on décrit les catégories $\mathrm{HOT}(S)$, conformément à la recette générale, à l'aide de \inserted{\unsure{contrav.}} foncteurs
\[
F:S^{\unsure{o|0|\circ}}\longrightarrow\mathrm{Cat}
\]
(préfaisceaux en catégories sur $S$), ou ce qui revient au $=$, de catégories scindées, la description de la \deleted{première} flèche $(*)$ est évidente : elle revient à la loi
\[
\left|\begin{aligned}
&F\longmapsto F\circ f^o\\
&\underline{\mathrm{Hom}}(S^o,\mathrm{Cat})\unsure{\text{---}|\longrightarrow}\underline{\mathrm{Hom}}(S'{}^o,\mathrm{Cat}),
\end{aligned}\right.
\]
en passant aux localisés \unsure{pour la|pour les} $W$-\unsure{équivalences|équivalence} argument par argument. Dans le langage des catégories fibrées, il est bien connu que cette opération triviale sur les préfaisceaux en \unsure{catégorie|catégories}, correspond à l'opération \uline{image inverse} de catégories sur $S$ -- laquelle opération passe aux scindages.
Cela nous donne \deleted{\unsure{donc}}
\[
f^*(\mathrm{hot}_{\deleted{\unsure{X}}W,S}(X))=\mathrm{hot}_{W,S'}(\underbrace{X\times_S S'}_{X'}),
\]
pour \deleted{\unsure{tout}} $X$ fibrée scindée sur $S$. Or on sait que le scindage \inserted{de $X$} ne joue aucun rôle, \unsure{finalement|justement}, dans la description de la classe
\[
\mathrm{hot}_{W,S}(X)\in\mathrm{HOT}_W(S)
\]
% archive sheet number: 170
\scanpage{172}\mspage{117}
associée à $X$. On peut oublier le scindage, et même (pour \unsure{calculer|étudier} les $\mathrm{Hom}_{\mathrm{HOT}_W(S)}(X,Y)$) travailler avec des $S$-foncteurs \uline{quelconques} entre catégories fibrées sur $S$ -- pas la peine même de les supposer cartésiens, voire $=$ compatibles avec des scindages. Cela me paraît une chose \unsure{fort|très} \unsure{importante} pour le formalisme des $\mathrm{HOT}_W(S)$, et très spéciale au cas \deleted{\unsure{des}} de la catégorie de modèles $\mathrm{Cat}$.
C'est \inserted{un peu} \unsure{tacite|traité} dans le théorème-scholie -- mais mes doutes initiaux concernant les catégories $\underline{\mathrm{Fibsc}}$ elles-mêmes montrent bien que je sentais qu'il y avait là quelque chose de pas évident, ni même \unsure{(pour moi)|/ pas gros|/ par gros} de vraisemblable. Et pourtant l'existence du foncteur $\unsure{\Phi|Q}$ montre qu'il en est bien ainsi.

On peut donc interpréter $f^*$ en termes de $\underline{\mathrm{Fib}}\,S$, $\underline{\mathrm{Fib}}\,S'$, par le foncteur chg$^{\mathrm t}$ de base
\[
f^*:\underline{\mathrm{Fib}}\,S\longrightarrow\underline{\mathrm{Fib}}\,S',\qquad X\longmapsto X\times_S S'\,/S',
\]
lequel foncteur préserve évidemment les $W$-équivalences par fibres $W_S^f$, $W_{S'}^f$, qui est celle \unsure{pertinentes} ici
\[
f^*(W_S^f)\subset W_{S'}^f\deleted{\unsure{\}}},\quad\unsure{d}
\]
de sorte que $f^*$ induit bien
\[
(\underline{\mathrm{Fib}}\,S)(W_S^f)^{-1}\longrightarrow(\underline{\mathrm{Fib}}\,S')(W_S^f)^{-1},
\]
% archive sheet number: 171
\scanpage{173}\mspage{118}
i.e. le foncteur $(*)$
\[
f^*:\mathrm{HOT}_W(S)\longrightarrow\mathrm{HOT}_W(S').
\]

D'autre part, $\mathrm{HOT}_W(S)$ peut être \unsure{décrit|décrite} aussi comme \unsure{localisé|localisée} de $\mathrm{Cat}_{/S}$ \deleted{\unsure{tout}} entier, par rapport au localiseur $W_S^d$, \deleted{qui}\inserted{lequel} sur $\underline{\mathrm{Fib}}\,S$, et $=$ sur $\underline{\mathrm{Liss}}\,S$, coïncide avec $W_S^f$. Mais on se bute ici à la difficulté que le foncteur image inverse
\[
f^*:\mathrm{Cat}_{/S}\longrightarrow\mathrm{Cat}_{/S'},\quad X\longmapsto X\times_S S'\,/S'
\]
n'est \uline{pas} compatible en général (si $S$ n'est \deleted{\unsure{équivalente}}\inserted{isomorphe} à la catégorie ponctuelle, pour $S'$ quelconque) avec les localiseurs $W_S^d,W_{S'}^d$. On le voit p.ex. en prenant $S'=\{s\}$, et en se \unsure{rappélant|rappelant} que $W_S^d$ n'implique \uline{pas}, dans $\mathrm{Cat}_{/S}$ tout entier, la $W$-équivalence par fibres.

Ces ennuis n'ont pas lieu sur $\underline{\mathrm{Liss}}\,S$, où $W_S^d$ et $W_S^f$ coïncident. On peut donc décrire $f^*$ \inserted{sur les $\mathrm{HOT}$} aussi à l'aide de
\[
f^*:\underline{\deleted{\unsure{\mathrm{Fib}}}\mathrm{Liss}}\,S\longrightarrow\underline{\mathrm{Liss}}\,S':\quad X\longrightarrow X\times_S S'\,/S',
\]
en passant aux localisés
\[
f^*:(\underline{\mathrm{Liss}}\,S)(W_S^f)^{-1}\longrightarrow(\underline{\mathrm{Liss}}\,S')(W_{S'}^f)^{-1}.
\]

En d'autres termes :

\uline{Prop 1} Si $X$ dans $\mathrm{Cat}_{/S}$ est \uline{lisse} sur $S$, alors on a un isom. can.
\[
f^*(\mathrm{hot}_{W,S}(X))\simeq\mathrm{hot}_{W,S'}(X'),\quad\text{où }X'=X\times_S S'\,/\unsure{S'}.
\]
% archive sheet number: 172
\scanpage{174}\mspage{119}
Une remarque similaire s'applique à $\underline{\unsure{\mathrm{Fibhot}|\mathrm{Fiblat}}}\,S$, où $W_S^d=W_S^f$ également. Le foncteur chgt de base induit
\[
\underline{\mathrm{Fibhot}}\,S\longrightarrow\underline{\mathrm{Fibhot}}\,S',\quad X\longmapsto X\times_S S'\,/S',
\]
compatible avec $\deleted{\unsure{W_S^d}\,\unsure{W_S^f}}$ les localiseurs relatifs pertinents (tous ceux envisagés jusqu'à présent coïncident sur ces catégories, avec les hyp. faites sur $W$), d'où
\[
f^*:(\underline{\mathrm{Fibhot}}\,S)W_S^{-1}\longrightarrow(\underline{\mathrm{Fibhot}}\,S')W_{S'}^{-1}.
\]

Les localisés cette fois ne sont pas \inserted{les} $\mathrm{HOT}_W$, mais les $\mathrm{HOT}^{\mathrm{lc}}_W$, donc on trouve une description sympathique de
\[
f^*:\mathrm{HOT}^{\mathrm{lc}}_W(S)\longrightarrow\mathrm{HOT}^{\mathrm{lc}}_W(S'),
\]
plus avantageuse \unsure{encore|même} qu'en \deleted{\unsure{fonct.}} termes de $\underline{\mathrm{Liss}}_0\,S$, $\underline{\mathrm{Liss}}_0\,S'$, car s'appliquant à des catégories \inserted{un peu} plus vastes (NB $\underline{\mathrm{Liss}}_0\subset\underline{\mathrm{Fibhot}}$).

Notons donc le

\uline{Corollaire 1} La formule de la prop.~1 \inserted{(p.~118)} reste valable si $X$ est supposé $\mathrm{Hot}$-fibrant sur $S$ (au lieu de lisse sur $S$). (On trouve alors des types d'homotopie relatifs \uline{loc. const.} sur $S,S'$.)

On a donc une description de $f^*$ qui s'applique à \unsure{toutes} les catégories du côté gauche
% archive sheet number: 173
\scanpage{175}\mspage{120}
du diagramme 76, à la seule exception de $\mathrm{Cat}_{/S}$ elle-même. Si $X$ est dans $\mathrm{Cat}_{/S}$, sans être ni lisse ni Hot-fibrée sur $S$, comment trouver \unsure{l'|son }image inverse dans $\mathrm{HOT}_W(S')$ ? On a le choix entre $\unsure{\Phi|Q}$ et $C$ pour améliorer la situation ! On sait que $X$ définit la $=$ chose dans $\mathrm{HOT}_W(S)$ que $\Phi(X),C(X)$, qui sont l'une lisse (et $=$ fibrée scindée), l'autre Hot-fibrée\deleted{\unsure{s}} (et même parfaite) sur $S$. Donc

\uline{Corollaire 2} Si $X$ dans $\mathrm{Cat}_{/S}$ est quelconque, alors
\[
f^*(\mathrm{hot}_{W,S}(X))\simeq\mathrm{hot}_{W,S'}(\Phi(X)\times_S S')\simeq\mathrm{hot}_{W,S'}(C(X)\times_S S').
\]

On peut s'attendre, d'autre part, que pour certains changements de base particuliers, la formule de la prop.~1 soit valable sans restriction sur $X$ dans $\mathrm{Cat}_{/S}$. Il doit revenir $\pm$ au même que le foncteur image inverse
\[
X\longmapsto X\times_S S'
\]
est compatible avec les localiseurs $W_S^d,W_{S'}^d$.

La \deleted{\unsure{con}} condition de $W$-équivalence ``colocale sur $S$'' doit être vue comme une condition de nature ``homologique'', pour le dérivateur $\mathrm{HOT}$. Elle s'associe, pour un dérivateur quelconque $\mathbb D$, à la condition (pour $f:X\longrightarrow Y$ de $\mathrm{Cat}_{/S}$) que \deleted{les} les $ {}_{s\backslash}X\unsure{\text{---}|\longrightarrow}{}_{s\backslash}Y$ soient des $\mathbb D$-équivalences,
% archive sheet number: 174
\scanpage{176}\mspage{121}
ce qui s'écrit \unsure{ici|très} élégamment en termes d'homologie relative, par la condition que $\forall a\in\unsure{\mathcal A|A}=\mathbb D(e)$, la flèche induite par $f$
\diagram{\begin{tikzcd}[ampersand replacement=\&,column sep=large,row sep=small]
H_{\unsure{\bullet|\cdot}}^{\mathbb D}(X/S,a_X) \arrow[r] \arrow[d,equal,"\mathrm{def}"'] \& H_{\bullet}^{\mathbb D}(\deleted{\unsure{X}}Y/S,a_Y) \arrow[d,equal,"\mathrm{def}"] \\
\deleted{\unsure{p}}{p_X}_{!}(a_X) \& (p_Y)_{!}(a_Y)
\end{tikzcd}}
soit iso. (Utiliser le ``calcul'' des faisceaux d'homologie relative $(p_X)_!(a_X)$, $(p_Y)_!(a_Y)$, par la caractérisation des fibres de Der 4 bis).

Or, les changements de base $\mathbb D$-\uline{propres} $S'\longrightarrow S$ commutent aux $p_!$, de sorte qu'on trouve le

\uline{Lemme} La notion de \deleted{\unsure{$W$}}-équivalence $\mathbb D$-colocale sur $S$ dans $\mathrm{Cat}_{/S}$ est stable par \inserted{\unsure{tout}} changement de base $S'\longrightarrow S$ $\mathbb D$-propre. (Quel que soit le dérivateur $\mathbb D$)

\deleted{Cela} On se rappellera que les $ {}_{s\backslash}S\longrightarrow S$ sont justement des changements de base propres ! Ceci \unsure{peut|veut} dire que l'équivalence $\mathbb D$-colocale peut se décrire en exigeant que $f:X\longrightarrow Y$ soit transformé en $W$-équivalence par certains types de chgt. de base $S'\longrightarrow S$, en prenant au choix l'un des trois types

a) les chgts de base $ {}_{s\backslash}S\longrightarrow S$

(ce qui redonne la définition que nous avions \unsure{prise|posée})

b) les chgts de base qui sont des co-isom. locaux

c) les chgt. de base propres

d) les chgts de base $\mathbb D$-propres.

On ne peut appliquer ces résultats au \unsure{dériv.-|dériva-}
% archive sheet number: 175
\scanpage{177}\mspage{122}
$\mathrm{HOT}$, faute de l'avoir encore construit, \deleted{\unsure{ou}}\inserted{et} vérifié les axiomes pertinents. On peut poser simplement la

\uline{Conjecture} (sûre !) : Soit $p:S'\longrightarrow S$ une flèche propre (sous-entendu : \unsure{Hot|HOT}-propre) dans $\mathrm{Cat}$. Alors le foncteur chgt de base
\[
\mathrm{Cat}_{/S}\unsure{\text{---}|\longrightarrow}\mathrm{Cat}_{/S'},\qquad X\longmapsto X\times_S S'\,/S'
\]
est compatible avec les localiseurs $W_S^d,W_S^{d'}$.
% slanted
\margin{NB Cette ``conjecture'' est prouvée sans bavure ! \unsure{le voir pp.} 168--169, sans avoir besoin de rien \unsure{de ce qui} ne soit déjà connu ici. Cf \unsure{th.~4, cor.~3)} p.~162}

(Pour un foncteur $S'\longrightarrow S$ lisse, le chang$^{\mathrm{mt}}$ de base serait compatible avec $W_S^g,W_{S'}^g$.)

Mais pour tt chgt de base $S'\longrightarrow S$ (propre ou non) qui a la propriété précédente, donc qui induit un foncteur
\diagram{\begin{tikzcd}[ampersand replacement=\&,column sep=large,row sep=small]
(\mathrm{Cat}_{/S})(W_S^d)^{-1} \arrow[r] \arrow[d,phantom,"{\unsure{\approx|\simeq}}"{description,sloped}] \& (\mathrm{Cat}_{/S'})(W_{S'}^d)^{-1} \arrow[d,phantom,"\approx"{description,sloped}] \\
\mathrm{HOT}_W(S) \& \mathrm{HOT}_W(S'),
\end{tikzcd}}
il est immédiat, à l'aide des th. de comparaison (p.ex. avec $\underline{\mathrm{Liss}}\,S$, $\underline{\mathrm{Liss}}\,S'$) sur $S$ et sur $S'$, que le foncteur n'est autre (à isom can. près) que le foncteur $f^*$ image inverse des types d'homotopie relatifs. \unsure{Donc à propos de ce|Donc à nouveau le}

\uline{Corollaire 3} : La formule de la prop.~\unsure{2|1} reste valable sans hypothèse aucune sur $X$ de $\mathrm{Cat}_{/S}$, pourvu que la flèche de chgt de base $f:S'\longrightarrow S$ soit telle que le foncteur $X\longmapsto X\times_S S'$ de $\mathrm{Cat}_{/S}$ \unsure{ds|de} $\mathrm{Cat}_{/S'}$
% archive sheet number: 176
\scanpage{178}\mspage{123}
% archive sheet number: 177
transforme $W$-\unsure{équivalences colocales|équivalence colocale} sur $S$ en $W$-\unsure{équivalences colocales|équivalence colocale} sur $S'$. (Ce qui devrait être le cas si (et seulement si ?) $f$ est \deleted{\unsure{$p$}} $W$-propre, donc aussi si $f$ est (``absolument'') propre\textsuperscript{*}.)
% slanted
\margin{cf I \unsure{p.~5, 52, 59|p.~5, 53, 59}}

On n'a travaillé jusqu'ici qu'avec les localiseurs $W_S^d,W_{S'}^d$, à l'\uline{exclusion} des $W_S,W_{S'}$ \inserted{(sauf sur les catégories, \unsure{convenablement} \uline{fibrées} /$S$, où \deleted{\unsure{elles}}les deux coïncident)}. Or, si on ne fait pas d'hypothèse draconienne sur $p:S'\longrightarrow S$, du type Hot-fibration, on n'a aucune \deleted{\unsure{raison}} chance pour que $X\longmapsto X\times_S S'$ transforme $W$-équivalences grossières dans itou, donc aucune chance d'attraper l'image inverse du type d'homotopie relatif ``grossier'' de $X$, en prenant celui de $X'=X\times_S S'$.

Ceci est lié au fait que pour un $S'\longrightarrow S$ général, les foncteurs de localisation
\[
\mathrm{HOT}_W(S)\xrightarrow{\rho_S}\mathrm{HOT}^{\mathrm{lc}}_W(S),\qquad
\mathrm{HOT}_W(S')\xrightarrow{\rho_{S'}}\mathrm{HOT}^{\mathrm{lc}}_W(S')
\]
n'ont aucune envie de commuter au foncteur image inverse $f^*$ sur les $\mathrm{HOT}_W$. On peut \unsure{formuler|transcrire} la

\uline{Proposition 2} Conditions équivalentes sur le foncteur $S'\unsure{\text{---}|\longrightarrow}S$ :

a) Le diagramme de foncteurs
\diagram{\begin{tikzcd}[ampersand replacement=\&]
\mathrm{HOT}_W(S)\arrow[r,"\rho_S"]\arrow[d,"f^*"] \& \mathrm{HOT}^{\mathrm{lc}}_W(S)\arrow[d,"f^*_{\mathrm{lc}}"]\\
\mathrm{HOT}_W(S')\arrow[r,no head,"\rho_{S'}"] \& \mathrm{HOT}^{\mathrm{lc}}_W(S')
\end{tikzcd}}
commute (à isom. \deleted{fonctoriel} de foncteurs \unsure{composés} \unsure{près})
% slanted
\margin{Pas prouvé, seulement que b, b', b'' impliquent \unsure{a}. Pour l'implication inverse, il faudrait savoir que toute flèche de $\unsure{\mathcal F|F}'$ qui devient iso dans $\mathcal F'W_{S'}^{-1}$, est dans $W_{S'}$ (saturation forte de $W_{S'}$ dans $\mathcal F'$), chose que je ne sais pas. (\unsure{Hélas|hélas} je crois \unsure{que ce doit|qu'on doit} être vrai)}
